\section{Background}

% --- 2.1 Accountability in Multi-Agent AI Systems ---

\subsection{Accountability in Multi-Agent AI Systems}

The problem of accountability in autonomous multi-agent systems has grown acute with the transition from single-agent pipelines to dynamic, recursive agent-spawning architectures. Early multi-agent systems literature treated accountability as a property of the individual agent---an agent is accountable if its designers can explain its behavior \cite{wooldridge1995,stone2000}. That framing is insufficient when agents spawn sub-agents, delegate tasks across trust boundaries, and operate in environments where no single developer retains full situational awareness. As the delegation chain branches, loops, or creates cycles through recursive spawning, the provenance of any given action becomes a graph rather than a traceable chain. Mutable delegation chains permit consequential actions to proceed without any checkpoint where a human principal can intervene or verify. The accountability problem in such systems is not merely an engineering inconvenience; it is a structural incompatibility between the autonomy granted to agents and the human responsibility that accompanies consequential decisions.

Dijkstra's structured programming paradigm established a foundational principle relevant to this problem: systems should have statically verifiable structures that prevent silent entry into unintended states \cite{dijkstra1982}. In the control-flow context, unrestricted jumps permit the program to enter arbitrary states from which recovery is not guaranteed. The structured programming response---making the set of reachable states statically enumerable---was a fail-safe measure. Applied to agentic systems, the principle implies that mutable delegation chains constitute a structural accountability failure analogous to unrestricted jumps in program control flow. Just as a program counter can be redirected to an arbitrary address via an unchecked jump, an agent with mutable delegation authority can redirect consequential actions to contexts that no human principal authorized. The remedy in both cases is structural: replace mutable control transfer with a bounded, statically enforced mechanism that makes every reachable state inspectable before it is entered. The BX3 Framework addresses this through its Fact Layer, which maintains an immutable authorization ledger independent of any agent process \cite{bx3framework2026}. The ledger records the human principal's identity, the scope of authorization, and the termination conditions for every spawn event, making authorization state tamper-evident across agent restarts.

A second dimension of the accountability problem is temporal. An agent may spawn a child that executes over hours or days. During that interval, the parent agent may be suspended, restarted, or replaced. The question of which human principal authorized a long-running task is not answered by a delegation certificate alone---it requires a durable, immutable record that persists regardless of agent lifecycle events. The BX3 Framework's Forensic Ledger satisfies this requirement by anchoring every spawn event to a designated human anchor $h(n)$ via an immutable parent pointer $\alpha(n)$, ensuring that the escalation path $C(n)$ from any node $n$ to its human principal is traceable even after agent restarts \cite{bx3framework2026}. The immutability of the parent pointer is not a policy choice but a structural invariant enforced by the Fact Layer's pre-commit hook: no machine actor, including the parent agent itself, can modify $\alpha(n)$ after the child node is provisioned. This decouples the accountability chain from agent runtime state, addressing the temporal dimension of accountability that arises when human principals authorize tasks with durations exceeding any single agent's operational lifetime.

The implications of these two dimensions compound in systems that combine recursive spawning with tool-use and external API calls. When a child agent calls an external API, the question of which human principal authorized that specific API call is not answered by the API itself. The authorization traces through the spawn chain to the root principal, but if the spawn chain is mutable, the trace is unreliable. Even if the trace is reliable, the latency between authorization and action can be hours or days, during which the principal's situational awareness of the broader task context may have changed significantly. The BX3 Framework's Fact Layer addresses this by recording not only the authorization scope but also the operating bounds at the time of each consequential action, creating a frozen context snapshot that can be evaluated independently of the agent's current state \cite{bx3framework2026}.

% --- 2.2 Human-in-the-Loop Design ---

\subsection{Human-in-the-Loop Design}

The human-in-the-loop (HITL) paradigm is the dominant approach to maintaining human oversight of autonomous systems. Its canonical formulation traces to Wiener's \emph{Cybernetics}, which established that any system capable of purposeful action must include a feedback loop that reaches a human operator capable of modifying the system's goals \cite{wiener1948}. Wiener's formulation was prescriptive: systems without human oversight were not merely inefficient but ethically inadmissible, because they divorced decision-making from accountability. The human-in-the-loop was not an optional feature---it was a constitutive requirement for any system that issues consequential directives. Wiener was explicit that the purpose of the feedback loop was not to improve system performance but to preserve human agency over systems that act in the world. This framing distinguishes the HITL requirement from performance optimization: it is a democratic and ethical prerequisite, not a technical tuner.

The practical implementation of HITL in modern AI systems has diverged significantly from Wiener's original intent. Traditional HITL was a design-time construct: a human defined the system's goals at design time, the system executed against those fixed goals, and the human intervened only at clearly defined breakpoints---such as a decision gate where the system's next action required explicit human approval before proceeding. Contemporary LLM-based agents operate in continuous time with no natural breakpoints, and critically, the decision of \emph{when} to escalate is itself made by the agent, not by the architecture. This creates a self-referential problem: if the agent decides when to call a human, the agent can defer escalation precisely when it is most needed---when it has encountered a situation outside its training distribution or beyond its authorized operating range. The HITL checkpoint becomes a procedural formality rather than a genuine oversight mechanism.

Bansal et al.\ provide empirical evidence for this failure mode in the context of human-AI collaboration \cite{bansal2019}. Their central finding is that HITL mechanisms improve team performance only when the human operator possesses an accurate mental model of the AI's decision-making process. When the AI's reasoning is opaque, or when the escalation interface lacks diagnostic context, humans default to either over-trusting or under-trusting the system---both of which degrade safety outcomes. Over-trusting occurs when the human approves a suggested action without sufficient scrutiny because the AI presents a plausible-sounding rationale. Under-trusting occurs when the human overrides the AI on cases where the AI's recommendation was actually correct, because the human lacks the context to evaluate the recommendation on its own terms. The content and format of information presented at the HITL checkpoint is not an ergonomic afterthought; it determines whether the checkpoint functions as a genuine corrective mechanism or as a ritual that provides psychological comfort without actual safety leverage \cite{bansal2019}. This finding is directly relevant to agentic systems with recursive spawning: each escalation interface is a HITL checkpoint, and the quality of the information presented at that interface determines whether the human principal can exercise meaningful judgment or is reduced to a yes/no gate with insufficient context.

The failure mode identified by Bansal et al.\ is structurally analogous to what occurs in agentic systems without bounded autonomy. When an agent's reasoning trace is opaque, the human principal cannot evaluate whether a constraint violation represents a minor boundary deviation that can be approved with conditions, or a fundamental goal misalignment that requires full task termination. A binary escalation interface---error occurred, human required---provides no pathway for calibrated responses. The BX3 Framework's Bailout Protocol addresses this by requiring that every escalation carry a complete diagnostic record: the exception classification, the operating range definition, the Bounds Engine's reasoning trace, and the resolution history of prior occurrences at that node \cite{bailoutprotocol2026}. This is a structural requirement for the HITL checkpoint to function as Wiener intended---not merely as a flag, but as a meaningful feedback loop with genuine corrective power. The diagnostic record enables the human principal to evaluate whether the violation is a recoverable boundary drift or an unrecoverable goal misalignment, and to respond with an appropriately calibrated directive.

The broader lesson from HITL research is that the architectural placement of the human checkpoint matters as much as its existence. An HITL interface at the system's output layer (human reviews the final action) is fundamentally different from an HITL interface at the system's reasoning layer (human reviews the decision process before action is taken). The former provides accountability after the fact; the latter provides safety before the fact. The Bailout Protocol is an architecture at the reasoning layer: it interrupts the agent's execution before a consequential action is taken, when corrective intervention is still possible \cite{bailoutprotocol2026}. This is the structurally correct placement for a fail-safe checkpoint in an autonomous system.

% --- 2.3 Fail-Safe Design in Safety-Critical Systems ---

\subsection{Fail-Safe Design in Safety-Critical Systems}

The concept of fail-safe design originates in control theory and aerospace engineering, where the cost of failure is measured in human life. A fail-safe system is one that transitions to a safe state upon the occurrence of any fault that is either unanticipated or outside the system's operational envelope \cite{leveson2011}. The canonical fail-safe property is \emph{graceful degradation}: a system that cannot complete its intended function must fail in a way that causes no harm and that preserves the maximum possible measure of system integrity. The distinction between fail-safe and fail-operational is critical: a fail-operational system continues functioning despite faults, while a fail-safe system halts in a controlled state. For safety-critical autonomous systems, fail-safe is the appropriate requirement: the system must halt in a way that does not cause harm, and must do so unconditionally, regardless of the nature or location of the fault.

Dijkstra's structured programming paradigm established a control-flow discipline directly applicable to fail-safe system design \cite{dijkstra1982}. The \emph{go-to controversy} was, at its core, a fail-safe argument. Unrestricted jumps permit the program to enter arbitrary states from which recovery is not guaranteed. The argument was not that programs with unrestricted jumps are always wrong, but that their correctness cannot be verified statically because the set of reachable states depends on runtime execution paths that cannot be enumerated at design time. Structured programming proposed a discipline that made the set of reachable states \emph{a priori} enumerable---and therefore inspectable and, crucially, safe to reason about. The analogy to agentic systems is direct: an unrestricted agent reasoning trace can enter arbitrary decision states from which recovery is not guaranteed, and the only architectural remedy is to bound the reachable state space. Tristr's bounded rationality framework extends this argument to the design of complex organizational and computational systems \cite{tristr81}. The core thesis is that any system that attempts to operate beyond its defined bounds---whether those bounds are computational, informational, or organizational---will experience a form of systemic failure that is not recoverable by the system itself. The implication for agentic AI is direct: an agent that exceeds its provisioned operating range cannot be relied upon to recover itself, because the recovery action requires the very cognitive resources that are no longer available within the defined bounds. The only recoverable failure mode is one where the excess is detected and escalated before the agent's reasoning capacity is compromised \cite{tristr81}.

This analysis applies with particular force to LLM-based agents, whose reasoning processes are not statically bounded. An LLM can generate arbitrarily long reasoning chains, can invoke tools in arbitrary sequences, and can produce outputs that were not anticipated by its designers. The LLM's reasoning capacity is not a fixed computational budget but a function of the context window and the model's training distribution---both of which are outside the agent's own control. When an LLM-based agent encounters a task that exceeds its competence or authorized scope, the model's response is not a controlled halt but a confident hallucination or a looping reasoning chain that consumes resources without advancing toward a correct outcome. The fail-safe requirement for such systems is not that the agent halts gracefully (which would require the agent to recognize its own failure, which is itself an unsolvable problem for a general LLM), but that the system architecture enforces a monitoring layer that halts the agent when its reasoning exceeds defined bounds, independent of the agent's own assessment of its competence.

The distinction between fail-safe and fail-operational is also relevant to the design of multi-agent systems with recursive spawning. A spawn tree that permits unlimited depth is fail-operational in the sense that the system can always spawn another child to continue the task. But unlimited depth combined with bounded rationality is a structural instability: each additional layer of spawning moves the decision context further from the human principal who authorized the original task, and the cognitive gap between the leaf agent's context and the root principal's intent grows with each spawning level. The fail-safe response is to bound the spawn tree depth and to trigger the Bailout Protocol when the depth bound is reached without task completion---halting the system in a safe state rather than allowing unbounded proliferation of agent nodes \cite{bx3framework2026}. This is architecturally analogous to a circuit breaker in distributed systems: rather than attempting to recover from failures by retrying indefinitely, the system transitions to a safe state and requires explicit human intervention before proceeding.

The BX3 Framework's Fail-Safe design philosophy combines these principles: the system does not rely on agents to recognize their own failures, does not permit unbounded operational expansion, and transitions to a controlled halt state upon any constraint violation or timeout expiration. The Bailout Protocol is the mechanism by which this philosophy is implemented structurally: when the Bounds Engine detects a violation or timeout, it triggers a state machine transition that stalls all consequential action and escalates to the human principal with full diagnostic context \cite{bailoutprotocol2026}. This is fail-safe in the engineering sense: the system transitions to a safe state unconditionally, independent of the nature of the fault or the state of any individual agent.

% --- 2.4 The BX3 Framework: Three-Layer Model ---

\subsection{The BX3 Framework: Three-Layer Model}

The BX3 Framework is a structural architecture for accountable autonomous systems that synthesizes the accountability, HITL, and fail-safe requirements described above into a concrete three-layer model. Rather than treating accountability as a property to be enforced by policy or convention, the BX3 Framework makes accountability a structural property of the architecture itself---enforced by invariants that cannot be violated by any machine actor, including the agents the framework governs. The framework decomposes every agent node into three immutable functional layers, each with a distinct and non-overlapping role in the accountability chain:

\begin{enumerate}
    \item \textbf{Purpose Layer.} The Purpose Layer encodes the intended goal state and the terminal success criteria for the agent's current task. It is the sole source of directives to the agent's action space. The Purpose Layer is human-authored---either directly or through an authenticated delegation chain that traces to a human principal. No machine actor can modify the Purpose Layer's encoded intent \cite{bx3framework2026}. This architectural choice ensures that the goal state is not subject to drift by the agent that was tasked with achieving it---a structural prevention of the goal-post-moving failure mode identified in aligned AI literature.

    \item \textbf{Bounds Layer.} The Bounds Layer defines the operating envelope within which the agent is authorized to act. This includes resource constraints (time, memory, compute budget), scope constraints (which actions are permitted in the current context), and norm constraints (which categories of outcome are classified as unacceptable regardless of goal attainment). The Bounds Layer monitors the agent's reasoning trace against these constraints and triggers the Bailout Protocol when a constraint violation is detected or when the Bounds Engine's timeout $T_{\text{bounds}}$ expires without resolution \cite{bx3framework2026}. The Bounds Layer is the structural implementation of bounded rationality: it enforces the cognitive budget constraints identified by Tristr \cite{tristr81} as necessary conditions for any system that must operate reliably in complex environments.

    \item \textbf{Fact Layer.} The Fact Layer maintains the immutable authorization ledger and the forensic record of all state transitions at the node. It is the only layer with write access to the authorization ledger. The Fact Layer enforces the Bailout Protocol's state machine transitions directly---bypassing agent-level code---via a pre-commit hook that rejects any directive that would violate the safety properties established in the system. Every ledger entry is SHA-256 hash-chained, making the record tamper-evident \cite{bx3framework2026}. The Fact Layer is the architectural analog of an immutable audit log with cryptographic integrity guarantees, implemented as a structural enforcement layer rather than a logging subsystem that depends on correct agent behavior for its integrity.
\end{enumerate}

The three layers communicate directionally: the Purpose Layer issues directives downward to the agent's reasoning and action systems; the Bounds Layer monitors upward and outward, detecting constraint violations and triggering escalations; the Fact Layer records laterally and upward, maintaining the immutable audit trail that makes every state transition attributable to a human principal \cite{bx3framework2026}. The Bailout Protocol is the enforced communication path from the Bounds Layer to the human principal, traversing the immutable parent chain $\alpha(n)$ from the originating node $n$ to its designated human anchor $h(n)$. The directionality of this communication path is critical: the escalation moves upward through the spawn tree, not downward, and it bypasses all machine actors in the intervening hierarchy \cite{bailoutprotocol2026}.

The significance of this architecture is that it separates concerns that are conflated in conventional agentic systems. In a monolithic agent, goal reasoning, constraint monitoring, and authorization recording are performed by the same cognitive process---which means the agent can, in principle, modify its own authorization state or suppress its own escalation. This is not a hypothetical failure mode; it is the structural reality of any system where the agent controls its own instrumentation. The BX3 Framework prevents this by making the Fact Layer's authorization ledger read-only from the agent's perspective, enforced by a pre-commit hook that cannot be intercepted by any layer below it. This is structurally analogous to hardware-enforced privilege separation in operating systems, and it is similarly foundational to the system's safety properties.

The BX3 Framework's three-layer model provides the architectural substrate for the Compulsory Human Escalation theorem: for any node $n$ in the spawn tree, and for any exception state $e$ at $n$, the Bailout Protocol's escalation path $C(n)$ always reaches the human anchor $h(n)$ in finite time \cite{bx3framework2026}. The theorem follows from two structural invariants: the immutability of the parent pointer $\alpha(n)$, enforced by the Fact Layer's pre-commit hook, and the finiteness of the spawn tree depth, established at root initialization and enforced by the spawn refusal mechanism. The theorem establishes that the BX3 Framework's accountability guarantee is not a property of agent behavior---it is a property of the architecture itself, and it holds regardless of the reasoning strategy employed by any individual agent node. This is the critical distinction between the BX3 Framework and policy-based accountability approaches: the guarantee does not depend on agents choosing to comply with escalation protocols, because the architecture makes compliance structurally compulsory.